\begin{houseviewbox}[核心结论]
固定CNY200,000采用ETF/资产层优先的基准稳健核心模型：20\%组合内受保障流动性、5\%交易现金、35\%高质量人民币利率债、30\%红利权益ETF和10\%非杠杆黄金ETF。直接股票和公募REITs初始配置均为0\%。截至本报告，ETF执行日资格包尚未运行；任何未通过或缺少数据的ETF额度均留在受保障流动性。
\end{houseviewbox}

\section{模型标的与条件性动作}

对本案固定CNY200,000本金，以下是明确的条件性模型目标：先建立防守层；仅在执行日资格包通过后建立利率债，并在六个月内分批建立红利权益和黄金。它们是风险预算中的指定标的、金额、条件和后备位置，不是个性化适当性结论、交易委托或以历史价格推导的数量。

\begin{exhibitbox}[固定CNY200,000的指定模型标的]
\small
\begin{tabularx}{\exhibitboxwidth}{L{2.4cm}L{3.0cm}L{1.7cm}P}
\toprule
\textbf{资产层} & \textbf{指定模型标的} & \textbf{目标金额} & \textbf{条件性模型动作} \\
\midrule
受保障流动性 & 合格受保障存款或期限匹配的直接政府债 & CNY40,000 & T0建立；不以基金或理财替代受保障层。 \\
交易现金 & 账户结算现金或银行活期；\texttt{511880}仅为可选补充 & CNY10,000 & T0保留，用于分批实施和再平衡；硬上限为CNY10,000。 \\
利率债 & \texttt{511010} 国泰上证5年期国债ETF & CNY70,000 & 仅在完整执行日资格包使$q=1$后于T0建立。 \\
红利权益A & \texttt{510880} 华泰柏瑞上证红利ETF & CNY30,000 & 仅在$q=1$后按六个月计划分批建立。 \\
红利权益B & \texttt{512890} 华泰柏瑞中证红利低波动ETF & CNY30,000 & 仅在$q=1$且与A的合并持仓穿透通过后分批建立。 \\
黄金 & \texttt{518880} 华安易富黄金ETF & CNY20,000 & 仅在$q=1$后按六个月计划分批建立；不使用期货或杠杆。 \\
REITs与直接股票 & 无初始标的 & CNY0 & 不买入；分别等待项目级或发行人级独立研究。 \\
\bottomrule
\end{tabularx}
\end{exhibitbox}

执行日资格包必须保存同一中国市场时段的来源、时间戳、产品文件、价格/IOPV、折溢价、买卖价差、流动性、账户和最小交易单位资料；对\texttt{512890}还须保存与\texttt{510880}合并后的持仓穿透。缺任一字段或未通过受控阈值时，相关$q=0$，不发生该笔交易。对应金额唯一留在合格受保障流动性；不会增加交易现金至CNY10,000以上，也不会自动转入\texttt{511010}、其他ETF、REITs或直接股票。当前资格包状态为未运行，因此本报告不表示任何ETF已通过。

\section{策略定位}

本案是一个长期、条件性的风险预算框架，服务于固定本金的流动性、分散和再平衡安排，不承诺某个分红率、票息、净值回报或现金提取额。

现金分配与总回报必须分开看待：ETF分配来自基金资产，除权日会影响净值；利率债ETF、货币工具和黄金ETF也不应被写成固定生活费来源。报告只记录实际税后、扣费后到账现金$D_0$，并按权重、上限和资格开关再投资。

\section{基准稳健核心版}

\begin{exhibitbox}[固定CNY200,000的风险预算]
\small
\begin{tabularx}{\exhibitboxwidth}{L{3.2cm}P P P}
\toprule
\textbf{资产层} & \textbf{权重} & \textbf{金额} & \textbf{角色与约束} \\
\midrule
组合内受保障流动性 & 20\% & CNY40,000 & 防守与用途匹配；仅合格存款或到期前直接政府债可归入该层，也承接未通过ETF核验的余额。 \\
交易现金 & 5\% & CNY10,000 & 分批实施与再平衡缓冲；不等同于存款保险，且不得超过CNY10,000。 \\
高质量人民币利率债 & 35\% & CNY70,000 & 利率与防守资产层；仍须接受久期、净值和流动性风险及执行日资格核验。 \\
红利权益ETF & 30\% & CNY60,000 & 条件性权益与红利因子暴露；\texttt{512890}未通过穿透时保持0。 \\
非杠杆黄金ETF & 10\% & CNY20,000 & 宏观风险分散；不计入收入预算。 \\
公募REITs & 0\% & CNY0 & 只有三只独立项目通过项目、估值和流动性审查时，才可从利率债桶转入最多5\%。 \\
直接高股息股票 & 0\% & CNY0 & 观察池，不是初始配置；未来仅可替换红利权益ETF暴露，合计不超过15\%。 \\
\midrule
\textbf{合计} & \textbf{100\%} & \textbf{CNY200,000} & \\
\bottomrule
\end{tabularx}
\end{exhibitbox}

\section{回报和损失的正确表述}

长期年化3\%--6\%是五年以上规划时的公式敏感性区间，不是“最可能收益”、年度收益率、产品分配率或收益保证。参数$r$仅表示税费、交易成本和基金费用之前的名义总回报算术输入；它不生成或估计实际现金$D_0$。基准稳健核心在12个月共同承压情景下的算术结果为-16.30\%，固定本金对应约CNY32,600；这用于检验风险承受与流动性，不是止损线、最大损失承诺或实际回撤触发。

\section{适用条件与直接股票边界}

{\small
\begin{itemize}
  \item 资金期限至少五年，且三年内没有确定的购房、教育、医疗、债务偿还或其他刚性用途；若用途改变，应另建资金计划。
  \item 模型标的须在执行日复核法律文件、费率、价格/IOPV、折溢价、价差、流动性、持仓重叠和账户适当性；本报告当前不含通过记录。
  \item 直接股票只有在独立完成价格、估值、派息覆盖、现金流、税务、行业集中度和流动性核验后，才能替换红利权益ETF暴露；合计上限15\%、单一发行人上限5\%。不使用杠杆、黄金期货或借款，分配下降也不自动触发加仓。
\end{itemize}
\par}
